\documentclass[12pt,a4paper]{article}
\usepackage[fontset=fandol]{ctex}
\usepackage{amsmath,amssymb,amsthm}
\usepackage{geometry}
\usepackage{braket}
\usepackage{hyperref}

\hypersetup{
    colorlinks=true,
    linkcolor=black,
    urlcolor=blue,
    bookmarksnumbered=true
}

\geometry{left=2.5cm,right=2.5cm,top=2.5cm,bottom=2.5cm}
\setcounter{tocdepth}{4}

\title{量子力学\\第一章：量子力学基本概念、理论与数学方法}
\author{}
\date{}

\begin{document}
\maketitle
\thispagestyle{empty}
\tableofcontents
\newpage

\section*{参考教材}
\addcontentsline{toc}{section}{参考教材}

\begin{enumerate}
    \item 《量子力学》(I, II)——曾谨言
    \item 《高等量子力学》——喀兴林
    \item \textit{Modern Quantum Mechanics} —— J. J. Sakurai
    \item \textit{Quantum Mechanics}（共三卷）—— C. Cohen-Tannoudji, B. Diu, F. Laloë
\end{enumerate}

\section{量子力学基本概念、理论与数学方法}

\subsection{基本原理}

\subsubsection{原理1：态矢量}

微观系统的状态由 \textbf{Hilbert 空间}中的矢量 $\ket{\psi(t)}$ 描述。

\begin{itemize}
    \item \textbf{Hilbert 空间}：定义了内积运算的完备复数线性空间。
\end{itemize}

\subsubsection{原理2：观测量与算符}

微观系统的力学量用 Hilbert 空间中的\textbf{厄米算符 (Hermitian Operator)} $\hat{A}$ 描述。

\begin{itemize}
    \item \textbf{本征值方程}：$\hat{A}\ket{a_i}=a_i\ket{a_i}$
    \item 力学量的可测值仅限于其本征值 $a_i$。
    \item \textbf{完备性}：态矢量可展开为 $\ket{\psi(t)}=\sum_i c_i(t)\ket{a_i}$。
    \item \textbf{概率诠释}：$|c_i|^2$ 给出测量得到 $a_i$ 的概率（需满足归一化 $\sum_i|c_i|^2=1$）。
    \item \textbf{波包坍缩 (Collapse)}：测量后得到 $a_i$，系统态变为对应的本征态 $\ket{a_i}$。
\end{itemize}

\subsubsection{原理3：对易关系（正则量子化）}

直角坐标下，坐标算符与动量算符的对易关系为
\[
[\hat{r}_i,\hat{p}_j]=i\hbar\delta_{ij},\qquad (i,j=x,y,z)
\]
其中 $\delta_{ij}$ 为 Kronecker delta 符号。
\[
[\hat{r}_i,\hat{r}_j]=[\hat{p}_i,\hat{p}_j]=0
\]

\subsubsection{原理4：时间演化}

态矢量随时间演化遵循\textbf{薛定谔方程}：
\[
i\hbar\frac{\partial}{\partial t}\ket{\psi(t)}=\hat{H}\ket{\psi(t)}
\]
其中 $\hat{H}$ 为系统的哈密顿算符。

\subsubsection{原理5：全同粒子}

全同粒子系统的波函数，对任意一对粒子交换具有对称性或反对称性：

\begin{itemize}
    \item \textbf{玻色子 (Boson)}：对称波函数（自旋为 $\hbar$ 的整数倍），服从 Bose--Einstein 统计。
    \item \textbf{费米子 (Fermion)}：反对称波函数（自旋为 $\hbar$ 的半奇数倍），服从 Fermi--Dirac 统计。
\end{itemize}

\subsection{Dirac 记号与常用算符}

\subsubsection{右矢 (Ket) 与左矢 (Bra)}

\begin{itemize}
    \item \textbf{右矢}：表示态矢量，如 $\ket{\psi(t)},\;\ket{\varphi},\;\ket{a_i},\;\ket{x},\;\ket{p},\;\ket{\ell,m}$。
    \item \textbf{左矢}：表示对偶矢量，如 $\bra{\psi(t)},\;\bra{\varphi},\;\bra{a_i}$。
    \item \textbf{对应关系}：$\ket{\psi}\leftrightarrow\bra{\psi}$；若 $\ket{\psi'}=\lambda\ket{\psi}$，则 $\bra{\psi'}=\lambda^*\bra{\psi}$。
\end{itemize}

\subsubsection{内积 (Inner Product)}

\begin{itemize}
    \item 定义：$\braket{\varphi|\psi}$（复数）。
    \item \textbf{性质}：$\braket{\varphi|\psi}=\braket{\psi|\varphi}^*$（复共轭）。$\braket{\psi|\psi}\ge 0$，且 $\braket{\psi|\psi}=1$ 时为归一化态。正交性：若 $\braket{\varphi|\psi}=0$，则两态正交。Schwarz 不等式：$|\braket{\varphi|\psi}|^2\le\braket{\varphi|\varphi}\braket{\psi|\psi}$。
\end{itemize}

\subsubsection{外积 (Outer Product)}

\begin{itemize}
    \item 形式：$\ket{\varphi}\bra{\psi}$，它是一个\textbf{算符}。
    \item \textbf{作用}：$(\ket{\varphi}\bra{\psi})\ket{\chi}=\braket{\psi|\chi}\ket{\varphi}$。
    \item \textbf{投影算符}：$\hat{P}_\varphi=\ket{\varphi}\bra{\varphi}$。作用：$\hat{P}_\varphi\ket{\chi}=\braket{\varphi|\chi}\ket{\varphi}$。
\end{itemize}

\subsubsection{线性与反线性算符}

\begin{itemize}
    \item \textbf{线性算符} $\hat{A}$：$\hat{A}(c_1\ket{\psi_1}+c_2\ket{\psi_2})=c_1\hat{A}\ket{\psi_1}+c_2\hat{A}\ket{\psi_2}$。
    \item \textbf{反线性算符} $\hat{B}$：$\hat{B}(c_1\ket{\psi_1}+c_2\ket{\psi_2})=c_1^*\hat{B}\ket{\psi_1}+c_2^*\hat{B}\ket{\psi_2}$。
\end{itemize}

\subsubsection{厄米共轭与厄米算符}

\begin{itemize}
    \item \textbf{厄米共轭} $\hat{A}^\dagger$：满足 $\bra{\varphi}\hat{A}^\dagger\ket{\psi}=\bra{\psi}\hat{A}\ket{\varphi}^*$。性质：$(\hat{A}\hat{B})^\dagger=\hat{B}^\dagger\hat{A}^\dagger$；$(\ket{\psi}\bra{\varphi})^\dagger=\ket{\varphi}\bra{\psi}$。
    \item \textbf{厄米算符 (Hermitian)}：$\hat{A}^\dagger=\hat{A}$。本征值为实数。不同本征值的本征态正交：$\braket{a_i|a_j}=\delta_{ij}$。任意算符可分解为厄米部分与反厄米部分之和：$\hat{A}=\frac{1}{2}(\hat{A}+\hat{A}^\dagger)+\frac{1}{2}(\hat{A}-\hat{A}^\dagger)$。
    \item \textbf{反厄米算符}：$\hat{A}^\dagger=-\hat{A}$。
\end{itemize}

\subsubsection{幺正算符 (Unitary Operator)}

\begin{itemize}
    \item \textbf{定义}：$\hat{U}^\dagger=\hat{U}^{-1}$。
    \item \textbf{性质}：保持内积不变。$\ket{\psi'}=\hat{U}\ket{\psi}$，$\braket{\psi'|\varphi'}=\braket{\psi|\varphi}$，保持矢量模长不变。
\end{itemize}

\subsubsection{反幺正算符 (Antiunitary Operator)}

\begin{itemize}
    \item \textbf{定义}：反线性且满足 $\bra{\hat{u}\psi|\hat{u}\varphi}=\braket{\varphi|\psi}=\braket{\psi|\varphi}^*$。
    \item 通常写作 $\hat{u}=\hat{U}\hat{K}$（$\hat{K}$ 为复共轭算符）。
\end{itemize}

\subsubsection{正规算符 (Normal Operator)}

\textbf{定义}：满足以下对易关系的算符 $\hat{N}$：
\[
[\hat{N},\hat{N}^\dagger]=0
\]

\textbf{常见特例}：
\begin{itemize}
    \item \textbf{厄米算符}：本征值 $\lambda$ 为实数。
    \item \textbf{反厄米算符}：本征值 $\lambda$ 为纯虚数。
    \item \textbf{幺正算符}：本征值 $\lambda$ 为模为 $1$ 的复数（$e^{i\varphi}$）。
\end{itemize}

\textbf{正交性}：
\[
\braket{\lambda_n|\lambda_m}=\delta_{nm}
\]

\textbf{谱展开 (Spectral Expansion)}：
\[
\hat{N}=\sum_n\lambda_n\ket{\lambda_n}\bra{\lambda_n}
\]

\subsubsection{\texorpdfstring{等距算符 (Isometric Operator) $\hat{V}$}{等距算符 (Isometric Operator)}}

\textbf{定义}：
\[
\hat{V}^\dagger\hat{V}=1,\qquad \text{但}\quad \hat{V}\hat{V}^\dagger\neq 1
\]

\textbf{性质分析}：

\begin{enumerate}
    \item \textbf{保内积（正向）}：设 $\hat{V}\ket{\psi}=\ket{\psi'}$，$\hat{V}\ket{\varphi}=\ket{\varphi'}$，则：
    \[
    \braket{\varphi'|\psi'}=\bra{\varphi}\hat{V}^\dagger\hat{V}\ket{\psi}=\braket{\varphi|\psi}
    \]
    内积保持不变。
    \item \textbf{不保内积（反向）}：设 $\hat{V}^\dagger\ket{\psi}=\ket{\tilde{\psi}}$，$\hat{V}^\dagger\ket{\varphi}=\ket{\tilde{\varphi}}$，则：
    \[
    \braket{\tilde{\varphi}|\tilde{\psi}}=\bra{\varphi}\hat{V}\hat{V}^\dagger\ket{\psi}\neq\braket{\varphi|\psi}
    \]
    内积发生改变。
\end{enumerate}

\par\smallskip
\noindent\textbf{注意：}等距变换不等同于幺正变换！等距变换不直接对应于物理态的变换（因为不是满射）。
\par\smallskip

\subsection{表象与表象变换}

\subsubsection{A-表象（离散谱）}

设厄米算符 $\hat{A}$（$\hat{A}^\dagger=\hat{A}$）具有分立谱：
\[
\hat{A}\ket{a_i}=a_i\ket{a_i}\qquad (i=1,2,\ldots,n)
\]
满足正交归一性：
\[
\braket{a_i|a_j}=\delta_{ij}
\]

\paragraph{完备性与态矢量展开}
任意态矢量 $\ket{\psi(t)}$ 可展开为：
\[
\ket{\psi(t)}=\sum_{i=1}^{n}C_i(t)\ket{a_i}
\]
代入左矢 $\bra{a_j}$ 可得展开系数：
\[
C_i(t)=\braket{a_i|\psi(t)}
\]
因此完备性关系可写为外积形式：
\[
\sum_{i=1}^{n}\ket{a_i}\bra{a_i}=1
\]

\paragraph{态矢量的矩阵表示}
在 $A$-表象中，态矢量 $\ket{\psi(t)}$ 表示为一个列矩阵（坐标）：
\[
\ket{\psi(t)}\doteq
\begin{pmatrix}
\braket{a_1|\psi(t)}\\
\braket{a_2|\psi(t)}\\
\vdots\\
\braket{a_n|\psi(t)}
\end{pmatrix}
=
\begin{pmatrix}
C_1(t)\\
C_2(t)\\
\vdots\\
C_n(t)
\end{pmatrix}
\]
其中 $|C_i(t)|^2$ 代表测量得到 $a_i$ 的概率。

左矢 $\bra{\psi(t)}$ 表示为行矩阵：
\[
\bra{\psi(t)}\doteq\big(\braket{a_1|\psi(t)}^*,\;\braket{a_2|\psi(t)}^*,\;\ldots,\;\braket{a_n|\psi(t)}^*\big)\equiv\psi^\dagger(t)
\]

\paragraph{本征矢的表示}
算符 $\hat{A}$ 的另一个本征态 $\ket{b_n}$ 在 $A$-表象中表示为：
\[
\ket{b_n}\doteq
\begin{pmatrix}
\braket{a_1|b_n}\\
\braket{a_2|b_n}\\
\vdots\\
\braket{a_n|b_n}
\end{pmatrix},\qquad
\bra{b_n}\doteq\big(\braket{a_1|b_n}^*,\;\ldots,\;\braket{a_n|b_n}^*\big)
\]

\paragraph{算符的矩阵表示}
算符 $\hat{F}$ 在 $A$-表象中的矩阵元定义为：
\[
F_{ij}=\bra{a_i}\hat{F}\ket{a_j}
\]
矩阵形式为：
\[
\hat{F}\doteq
\begin{pmatrix}
F_{11} & F_{12} & \cdots & F_{1n}\\
F_{21} & F_{22} & \cdots & F_{2n}\\
\vdots & \vdots & \ddots & \vdots\\
F_{n1} & F_{n2} & \cdots & F_{nn}
\end{pmatrix}
\equiv F
\]

\textbf{厄米共轭的矩阵关系}：
\[
F_{ji}^*=(F^\dagger)_{ij}\quad\Rightarrow\quad F^\dagger=(F^T)^*
\]
即厄米共轭等于转置加复共轭。

\textbf{特殊算符的矩阵形式}：
\begin{itemize}
    \item \textbf{厄米算符}：$F^\dagger=F$（若矩阵元为实数，则为实对称矩阵）。
    \item \textbf{幺正算符}：$F^\dagger=F^{-1}$（若矩阵元为实数，则为正交矩阵）。
    \item \textbf{$\hat{A}$ 自身的矩阵}：由于 $\hat{A}\ket{a_i}=a_i\ket{a_i}$，
    \[
    A_{ij}=\bra{a_i}\hat{A}\ket{a_j}=a_j\braket{a_i|a_j}=a_j\delta_{ij}
    \]
    故 $\hat{A}$ 在其自身表象中是对角矩阵：
    \[
    A=
    \begin{pmatrix}
    a_1 & & & \\
     & a_2 & & \\
     & & \ddots & \\
     & & & a_n
    \end{pmatrix}
    \]
    其谱表示为 $\hat{A}=\sum_i a_i\ket{a_i}\bra{a_i}$。
\end{itemize}

\paragraph{外积的矩阵表示}
外积 $\ket{\varphi}\bra{\psi}$ 对应的矩阵 $\tilde{F}$ 为：
\[
\tilde{F}_{ij}=\braket{a_i|\varphi}\braket{\psi|a_j}=\braket{a_i|\varphi}\braket{a_j|\psi}^*
\]
矩阵形式为列向量乘以行向量（外积运算）：
\[
\tilde{F}=
\begin{pmatrix}
\braket{a_1|\varphi}\\
\vdots\\
\braket{a_n|\varphi}
\end{pmatrix}
\big(\braket{a_1|\psi}^*,\;\cdots,\;\braket{a_n|\psi}^*\big)
=\varphi\,\psi^\dagger
\]

\paragraph{平均值的矩阵表示}
算符 $\hat{F}$ 在态 $\ket{\psi(t)}$ 中的平均值：
\[
\langle\hat{F}\rangle=\bra{\psi(t)}\hat{F}\ket{\psi(t)}=\psi^\dagger(t)\,F\,\psi(t)
\]

\paragraph{本征值方程的矩阵形式}
求解 $\hat{F}\ket{f}=f\ket{f}$ 等价于求解矩阵本征值方程：
\[
F\mathbf{f}=f\mathbf{f}
\]
其中 $\mathbf{f}=\begin{pmatrix}\braket{a_1|f}\\ \vdots\\ \braket{a_n|f}\end{pmatrix}$。

存在非平庸解的条件是久期方程 (Secular Equation) 成立：
\[
\det(F-f\mathbf{1})=0
\]
即：
\[
\begin{vmatrix}
F_{11}-f & F_{12} & \cdots & F_{1n}\\
F_{21} & F_{22}-f & \cdots & F_{2n}\\
\vdots & \vdots & \ddots & \vdots\\
F_{n1} & F_{n2} & \cdots & F_{nn}-f
\end{vmatrix}=0
\]

\subsubsection{坐标表象与动量表象（连续谱）}

\paragraph{坐标表象（一维）}

\textbf{本征值方程}：
\[
\hat{x}\ket{x}=x\ket{x}
\]

\textbf{正交性}：
\[
\braket{x'|x}=\delta(x-x')
\]

\textit{推导简述}：利用 $\bra{x'}\hat{x}\ket{x}=x\braket{x'|x}=x'\braket{x'|x}$，得 $(x-x')\braket{x'|x}=0$。结合 $x\delta(x)=0$，得出比例关系，取归一化系数即得。

\textbf{完备性}：
\[
\int_{-\infty}^{\infty}dx\,\ket{x}\bra{x}=1
\]

\textbf{波函数表示}：
\[
\ket{\psi(t)}=\int_{-\infty}^{\infty}dx\,\ket{x}\braket{x|\psi(t)}=\int_{-\infty}^{\infty}dx\,\psi(x,t)\ket{x}
\]
其中 $\psi(x,t)=\braket{x|\psi(t)}$ 即为坐标空间波函数。

$|\psi(x,t)|^2$ 表示 $t$ 时刻在 $x$ 处找到粒子的概率密度。

\textbf{归一化条件}：
\[
\braket{\psi(t)|\psi(t)}=\int_{-\infty}^{\infty}dx\,\psi^*(x,t)\psi(x,t)=\int_{-\infty}^{\infty}dx\,|\psi(x,t)|^2=1
\]

\textbf{内积表示}：
\[
\braket{\varphi(t)|\psi(t)}=\int_{-\infty}^{\infty}dx\,\varphi^*(x,t)\psi(x,t)
\]

\paragraph{\texorpdfstring{Dirac $\delta$-函数性质}{Dirac delta-函数性质}}

\begin{enumerate}
    \item 奇偶性：$\delta(x)=\delta(-x)$
    \item 乘积性质：$x\delta(x)=0$
    \item 微分性质：$x\dfrac{d}{dx}\delta(x)=-\delta(x)$
    \item 积分定义：$\int_{-\infty}^{\infty}dx\,\delta(x)=1$（区间包含原点）
    \item 筛选性质：$\int_{-\infty}^{\infty}dx\,f(x)\delta(x)=f(0)$
    \item 复合性质：$\int_{-\infty}^{\infty}dx\,\delta(x-a)\delta(x-b)=\delta(a-b)$
    \item 缩放性质：$\delta(ax)=\dfrac{1}{|a|}\delta(x)$
\end{enumerate}

\textbf{初等函数表示}：

\begin{enumerate}
    \item[(a)] Lorentz 型：$\delta(x)=\lim_{\epsilon\to 0}\dfrac{1}{\pi}\dfrac{\epsilon}{x^2+\epsilon^2}$
    \item[(b)] 积分型：$\delta(x)=\dfrac{1}{2\pi}\int_{-\infty}^{\infty}dk\,e^{ikx}$
    \item[(c)] Gauss 型：$\delta(x)=\lim_{\epsilon\to 0}\dfrac{1}{\sqrt{\pi\epsilon}}\exp\!\left(-\dfrac{x^2}{\epsilon}\right)$
\end{enumerate}

\paragraph{坐标算符的矩阵表示}

坐标算符 $\hat{x}^n$ 的矩阵元：
\[
\bra{x}\hat{x}^n\ket{x'}=x^n\braket{x|x'}=x^n\delta(x-x')
\]
其谱表示为：
\[
\hat{x}^n=\int dx\,\ket{x}\,x^n\bra{x}
\]

\paragraph{坐标表象中动量算符的矩阵元}

动量算符 $\hat{p}$ 在坐标表象 $\{\ket{x}\}$ 下的矩阵元定义为：
\[
P_{xx'}=\bra{x}\hat{p}\ket{x'}=-i\hbar\frac{\partial}{\partial x}\delta(x-x')
\]

\textbf{证明}：

利用基本对易关系 $[\hat{x},\hat{p}]=i\hbar$，左乘 $\bra{x}$，右乘 $\ket{x'}$：
\[
\bra{x}[\hat{x},\hat{p}]\ket{x'}=i\hbar\braket{x|x'}
\]
左侧展开：
\[
\bra{x}(\hat{x}\hat{p}-\hat{p}\hat{x})\ket{x'}=(x-x')\bra{x}\hat{p}\ket{x'}
\]
右侧为：
\[
i\hbar\,\delta(x-x')
\]
因此：
\[
(x-x')\bra{x}\hat{p}\ket{x'}=i\hbar\,\delta(x-x')
\]
结合分布导数性质 $(x-x')\dfrac{\partial}{\partial x}\delta(x-x')=-\delta(x-x')$，可得：
\[
\bra{x}\hat{p}\ket{x'}=-i\hbar\frac{\partial}{\partial x}\delta(x-x')
\]

\paragraph{\texorpdfstring{坐标表象中 $\hat{p}^2$ 与 $\hat{p}^n$ 的矩阵元}{坐标表象中 p^2 与 p^n 的矩阵元}}

\textbf{(a) $\hat{p}^2$ 的矩阵元}：
\[
P_{xx'}^{(2)}=\bra{x}\hat{p}^2\ket{x'}
\]
插入单位算符 $\int_{-\infty}^{\infty}dx''\,\ket{x''}\bra{x''}$：
\begin{align*}
P_{xx'}^{(2)}&=\int_{-\infty}^{\infty}dx''\,\bra{x}\hat{p}\ket{x''}\bra{x''}\hat{p}\ket{x'}\\
&=\int_{-\infty}^{\infty}dx''\left[-i\hbar\frac{\partial}{\partial x}\delta(x-x'')\right]\left[-i\hbar\frac{\partial}{\partial x''}\delta(x''-x')\right]\\
&=\left(-i\hbar\frac{\partial}{\partial x}\right)\int_{-\infty}^{\infty}dx''\,\delta(x-x'')\left(-i\hbar\frac{\partial}{\partial x''}\delta(x''-x')\right)\\
&=\left(-i\hbar\frac{\partial}{\partial x}\right)\left(-i\hbar\frac{\partial}{\partial x'}\delta(x-x')\right)\\
&=\left(-i\hbar\frac{\partial}{\partial x}\right)^2\delta(x-x')
\end{align*}

\textbf{(b) $\hat{p}^n$ 的矩阵元}：

推广至 $n$ 次幂：
\[
P_{xx'}^{(n)}=\bra{x}\hat{p}^n\ket{x'}=\left(-i\hbar\frac{\partial}{\partial x}\right)^n\delta(x-x')
\]

\paragraph{坐标与动量算符混合期望值}

\textbf{(a) $\langle\hat{x}^n\hat{p}^m\rangle$}
\begin{align*}
\langle\hat{x}^n\hat{p}^m\rangle&=\bra{\psi}\hat{x}^n\hat{p}^m\ket{\psi}\\
&=\int dx_1\,dx_2\,dx_3\,\braket{\psi|x_1}\bra{x_1}\hat{x}^n\ket{x_2}\bra{x_2}\hat{p}^m\ket{x_3}\braket{x_3|\psi}\\
&=\int dx_1\,dx_2\,dx_3\,\psi^*(x_1)\,x_2^n\,\delta(x_1-x_2)\left(-i\hbar\frac{\partial}{\partial x_2}\right)^m\delta(x_2-x_3)\,\psi(x_3)\\
&=\int dx_2\,\psi^*(x_2)\,x_2^n\left(-i\hbar\frac{\partial}{\partial x_2}\right)^m\int dx_3\,\delta(x_2-x_3)\psi(x_3)\\
&=\int dx\,\psi^*(x)\,x^n\left(-i\hbar\frac{\partial}{\partial x}\right)^m\psi(x)
\end{align*}

\textbf{(b) $\langle\hat{p}^m\hat{x}^n\rangle$}
\[
\langle\hat{p}^m\hat{x}^n\rangle=\bra{\psi}\hat{p}^m\hat{x}^n\ket{\psi}=\int dx\,\psi^*(x)\left(-i\hbar\frac{\partial}{\partial x}\right)^m\!\big[x^n\psi(x)\big]
\]

\textbf{(c) 对易关系 $[\hat{x},\hat{p}]=i\hbar$ 的坐标表象验证}

作用于任意态 $\ket{\psi}$：
\[
[\hat{x},\hat{p}]\ket{\psi}=i\hbar\ket{\psi}
\]
左乘 $\bra{x}$ 并插入单位算符：
\begin{align*}
&\int dx_1\,dx_2\,\bra{x}\hat{x}\ket{x_1}\bra{x_1}\hat{p}\ket{x_2}\braket{x_2|\psi}-\int dx_1\,dx_2\,\bra{x}\hat{p}\ket{x_1}\bra{x_1}\hat{x}\ket{x_2}\braket{x_2|\psi}\\
&=\int dx_1\,dx_2\,x\,\delta(x-x_1)\left(-i\hbar\frac{\partial}{\partial x_1}\delta(x_1-x_2)\right)\psi(x_2)\\
&\quad-\int dx_1\,dx_2\left(-i\hbar\frac{\partial}{\partial x}\delta(x-x_1)\right)x_2\,\delta(x_1-x_2)\psi(x_2)\\
&=x\left(-i\hbar\frac{\partial}{\partial x}\right)\psi(x)-\left(-i\hbar\frac{\partial}{\partial x}\right)\!\big[x\,\psi(x)\big]\\
&=i\hbar\,\psi(x)
\end{align*}
这正是坐标表象下的算符对易关系：
\[
\left[x,\,-i\hbar\frac{\partial}{\partial x}\right]\psi(x)=i\hbar\,\psi(x)
\]

\paragraph{能量本征态的正交完备性（坐标表象）}

设 $\hat{A}\ket{a_i}=a_i\ket{a_i}$，定义坐标空间波函数 $U_i(x)=\braket{x|a_i}$。

\textbf{(1) 正交性}：

由 $\braket{a_j|a_i}=\delta_{ij}$，插入单位算符 $\int dx\,\ket{x}\bra{x}$：
\[
\int dx\,\braket{a_j|x}\braket{x|a_i}=\int dx\,U_j^*(x)\,U_i(x)=\delta_{ij}
\]

\textbf{(2) 完备性}：

由 $\sum_i\ket{a_i}\bra{a_i}=1$，左乘 $\bra{x'}$，右乘 $\ket{x}$：
\[
\sum_i\braket{x'|a_i}\braket{a_i|x}=\braket{x'|x}=\delta(x'-x)
\]
即：
\[
\sum_i U_i^*(x')\,U_i(x)=\delta(x'-x)
\]

\paragraph{一维谐振子的坐标表象薛定谔方程}

哈密顿量：$\hat{H}=\dfrac{\hat{p}^2}{2m}+\dfrac{1}{2}m\omega^2\hat{x}^2$。

薛定谔方程 $i\hbar\dfrac{\partial}{\partial t}\ket{\psi(t)}=\hat{H}\ket{\psi(t)}$ 左乘 $\bra{x}$：
\[
i\hbar\frac{\partial}{\partial t}\psi(x,t)=\bra{x}\hat{H}\ket{\psi(t)}
\]
其中 $\psi(x,t)=\braket{x|\psi(t)}$。

\textbf{动能项}：
\begin{align*}
\bra{x}\frac{\hat{p}^2}{2m}\ket{\psi(t)}&=\frac{1}{2m}\int dx'\,\bra{x}\hat{p}^2\ket{x'}\braket{x'|\psi(t)}\\
&=\frac{1}{2m}\int dx'\left(-i\hbar\frac{\partial}{\partial x}\right)^2\delta(x-x')\,\psi(x',t)\\
&=\frac{1}{2m}\left(-i\hbar\frac{\partial}{\partial x}\right)^2\psi(x,t)\\
&=-\frac{\hbar^2}{2m}\frac{\partial^2}{\partial x^2}\psi(x,t)
\end{align*}

\textbf{势能项}：
\[
\bra{x}\frac{1}{2}m\omega^2\hat{x}^2\ket{\psi(t)}=\frac{1}{2}m\omega^2 x^2\psi(x,t)
\]

\textbf{最终方程}：
\[
i\hbar\frac{\partial}{\partial t}\psi(x,t)=\left[-\frac{\hbar^2}{2m}\frac{\partial^2}{\partial x^2}+\frac{1}{2}m\omega^2 x^2\right]\psi(x,t)
\]

\paragraph{能量本征值方程的坐标表象}

定态方程 $\hat{H}\ket{E}=E\ket{E}$，左乘 $\bra{x}$ 得：
\[
\left[-\frac{\hbar^2}{2m}\frac{\partial^2}{\partial x^2}+\frac{1}{2}m\omega^2 x^2\right]\psi_E(x)=E\,\psi_E(x)
\]
其中 $\psi_E(x)=\braket{x|E}$。

\textbf{解析解}：
\[
\begin{cases}
E_n=\left(n+\dfrac{1}{2}\right)\hbar\omega,\\[6pt]
\psi_n(x)=\braket{x|n}=\left(\dfrac{m\omega}{\pi\hbar}\right)^{1/4}\dfrac{1}{\sqrt{2^n n!}}\,H_n(\xi)\,e^{-\xi^2/2},\\[6pt]
n=0,1,2,\ldots
\end{cases}
\]
其中 $\xi=x/x_0$，$x_0=\sqrt{\hbar/(m\omega)}$，$H_n(\xi)$ 为 Hermite 多项式：
\[
H_n(\xi)=(-1)^n e^{\xi^2}\left(\frac{d}{d\xi}\right)^n e^{-\xi^2}
\]

\paragraph{动量表象}

\textbf{(a) 基本性质}

动量本征态 $\ket{p}$ 满足 $\hat{p}\ket{p}=p\ket{p}$。

\begin{itemize}
    \item \textbf{正交性}：$\braket{p'|p}=\delta(p'-p)$
    \item \textbf{完备性}：$\int_{-\infty}^{\infty}dp\,\ket{p}\bra{p}=1$
\end{itemize}

态矢量在动量表象的展开：
\[
\ket{\psi(t)}=\int_{-\infty}^{\infty}dp\,\ket{p}\braket{p|\psi(t)}=\int_{-\infty}^{\infty}dp\,\ket{p}\,\psi(p,t)
\]
其中 $\psi(p,t)=\braket{p|\psi(t)}$ 为动量空间波函数，$|\psi(p,t)|^2$ 表示动量为 $p$ 的概率密度。

\textbf{(b) 动量表象中的算符矩阵元}

\begin{itemize}
    \item \textbf{坐标算符}：$x_{pp'}=\bra{p}\hat{x}\ket{p'}=i\hbar\dfrac{\partial}{\partial p}\delta(p-p')$
    \item \textbf{坐标算符 $n$ 次幂}：$x_{pp'}^{(n)}=\bra{p}\hat{x}^n\ket{p'}=\left(i\hbar\dfrac{\partial}{\partial p}\right)^n\delta(p-p')$
    \item \textbf{动量算符 $n$ 次幂}：$P_{pp'}^{(n)}=\bra{p}\hat{p}^n\ket{p'}=p^n\,\delta(p-p')$
\end{itemize}

\textbf{(c) 坐标与动量表象的波函数关系}

由 $\braket{x|p}=\dfrac{1}{\sqrt{2\pi\hbar}}e^{ipx/\hbar}$，可得 Fourier 变换关系：
\begin{align*}
\psi(x,t)&=\frac{1}{\sqrt{2\pi\hbar}}\int_{-\infty}^{\infty}dp\,e^{ipx/\hbar}\,\psi(p,t)\\
\psi(p,t)&=\frac{1}{\sqrt{2\pi\hbar}}\int_{-\infty}^{\infty}dx\,e^{-ipx/\hbar}\,\psi(x,t)
\end{align*}

\paragraph{坐标与动量表象的变换关系}

算符 $\hat{F}$ 的矩阵元变换：
\begin{align*}
F_{xx'}&=\bra{x}\hat{F}\ket{x'}=\int dp\,dp'\,\braket{x|p}\bra{p}\hat{F}\ket{p'}\braket{p'|x'}\\
&=\frac{1}{2\pi\hbar}\int dp\,dp'\,e^{i(px-p'x')/\hbar}\,F_{pp'}\\[6pt]
F_{pp'}&=\bra{p}\hat{F}\ket{p'}=\int dx\,dx'\,\braket{p|x}\bra{x}\hat{F}\ket{x'}\braket{x'|p'}\\
&=\frac{1}{2\pi\hbar}\int dx\,dx'\,e^{-i(px-p'x')/\hbar}\,F_{xx'}
\end{align*}

\textbf{示例：求 $x_{pp'}$}

已知 $x_{xx'}=x\,\delta(x-x')$，代入变换公式：
\begin{align*}
x_{pp'}&=\frac{1}{2\pi\hbar}\int dx\,dx'\,e^{-i(px-p'x')/\hbar}\,x\,\delta(x-x')\\
&=\frac{1}{2\pi\hbar}\int dx\,x\,e^{-i(p-p')x/\hbar}\\
&=i\hbar\frac{\partial}{\partial p}\left[\frac{1}{2\pi\hbar}\int dx\,e^{-i(p-p')x/\hbar}\right]\\
&=i\hbar\frac{\partial}{\partial p}\delta(p-p')
\end{align*}

\subsection{一维线性谐振子的代数解法}

\subsubsection{产生与湮灭算符的定义}

引入无量纲算符：
\[
\begin{cases}
\hat{a}=\sqrt{\dfrac{m\omega}{2\hbar}}\left(\hat{x}+\dfrac{i}{m\omega}\hat{p}\right)\\[8pt]
\hat{a}^\dagger=\sqrt{\dfrac{m\omega}{2\hbar}}\left(\hat{x}-\dfrac{i}{m\omega}\hat{p}\right)
\end{cases}
\]
$\hat{a}$ 与 $\hat{a}^\dagger$ 互为厄米共轭，\textbf{不是厄米算符}。

\textbf{逆变换}：
\[
\begin{cases}
\hat{x}=\sqrt{\dfrac{\hbar}{2m\omega}}\left(\hat{a}+\hat{a}^\dagger\right)\\[8pt]
\hat{p}=i\sqrt{\dfrac{m\hbar\omega}{2}}\left(\hat{a}^\dagger-\hat{a}\right)
\end{cases}
\]

\subsubsection{对易关系}

利用 $[\hat{x},\hat{p}]=i\hbar$，可得：
\[
[\hat{a},\hat{a}^\dagger]=1
\]
且 $\hat{N}=\hat{a}^\dagger\hat{a}$ 为厄米算符（粒子数算符）。

\subsubsection{\texorpdfstring{粒子数算符 $\hat{N}$ 的本征值与本征态}{粒子数算符的本征值与本征态}}

本征方程：$\hat{N}\ket{n}=n\ket{n}$，$\braket{n|n}=1$。

\textbf{(a) 本征值 $n$ 的非负性}

由 $\bra{n}\hat{N}\ket{n}=n\braket{n|n}=n$，且 $\bra{n}\hat{a}^\dagger\hat{a}\ket{n}=\|\hat{a}\ket{n}\|^2\ge 0$，得 $n\ge 0$。

\textbf{(b) 递推关系与本征值谱}

利用对易关系可证：
\[
\hat{N}(\hat{a}^\dagger)^m=(\hat{a}^\dagger)^m(\hat{N}+m)
\]
作用于基态 $\ket{0}$（满足 $\hat{a}\ket{0}=0$）：
\[
\hat{N}(\hat{a}^\dagger)^m\ket{0}=m\,(\hat{a}^\dagger)^m\ket{0}
\]
表明 $(\hat{a}^\dagger)^m\ket{0}$ 是 $\hat{N}$ 本征值为 $m$ 的本征态。进一步由 $\bra{n}(\hat{a}^\dagger)^m\hat{a}^m\ket{n}=n(n-1)\cdots(n-m+1)\ge 0$，推得 $n$ 必须为非负整数：
\[
n=0,1,2,\ldots
\]

\textbf{(c) 归一化本征态}
\[
\ket{n}=\frac{1}{\sqrt{n!}}(\hat{a}^\dagger)^n\ket{0}
\]
满足正交归一性 $\braket{n'|n}=\delta_{n'n}$ 和完备性 $\sum_{n=0}^{\infty}\ket{n}\bra{n}=1$。

\subsubsection{产生与湮灭算符的作用}
\begin{align*}
\hat{a}^\dagger\ket{n}&=\sqrt{n+1}\,\ket{n+1} \quad &\text{（上升算符）}\\
\hat{a}\ket{n}&=\sqrt{n}\,\ket{n-1} \quad &\text{（下降算符，$n>0$）}\\
\hat{a}\ket{0}&=0
\end{align*}

\subsubsection{哈密顿量与能谱}
\[
\hat{H}=\hbar\omega\left(\hat{a}^\dagger\hat{a}+\frac{1}{2}\right)=\hbar\omega\left(\hat{N}+\frac{1}{2}\right)
\]
本征能量：
\[
E_n=\left(n+\frac{1}{2}\right)\hbar\omega,\qquad n=0,1,2,\ldots
\]

\begin{itemize}
    \item \textbf{基态}：$\ket{0}$，$E_0=\dfrac{1}{2}\hbar\omega$
    \item \textbf{激发态}：$\ket{n}$，$E_n-E_0=n\hbar\omega$
\end{itemize}

\subsubsection{坐标表象下的谐振子波函数}

\textbf{(a) 基态波函数 $\psi_0(x)=\braket{x|0}$}

由 $\hat{a}\ket{0}=0$，代入坐标表象：
\[
\sqrt{\frac{m\omega}{2\hbar}}\left(x+\frac{\hbar}{m\omega}\frac{\partial}{\partial x}\right)\psi_0(x)=0
\]
化简为：
\[
\left(x+\frac{\hbar}{m\omega}\frac{\partial}{\partial x}\right)\psi_0(x)=0
\quad\Rightarrow\quad
\frac{d\psi_0}{dx}=-\frac{m\omega}{\hbar}x\,\psi_0
\]
解得：
\[
\psi_0(x)=A\exp\!\left(-\frac{m\omega}{2\hbar}x^2\right)=A\exp\!\left(-\frac{x^2}{2x_0^2}\right)
\]
由归一化 $\int_{-\infty}^{\infty}|\psi_0(x)|^2\,dx=1$ 得：
\[
A=\left(\frac{m\omega}{\pi\hbar}\right)^{1/4}=\frac{1}{\pi^{1/4}\sqrt{x_0}}
\]

\textbf{(b) 激发态波函数}

利用 $\bra{x}\hat{a}^\dagger=\sqrt{\dfrac{m\omega}{2\hbar}}\left(x-\dfrac{\hbar}{m\omega}\dfrac{\partial}{\partial x}\right)\bra{x}$ 和 $\ket{n}=\dfrac{1}{\sqrt{n!}}(\hat{a}^\dagger)^n\ket{0}$，得：
\begin{align*}
\psi_n(x)&=\braket{x|n}=\frac{1}{\sqrt{n!}}\bra{x}(\hat{a}^\dagger)^n\ket{0}\\
&=\frac{1}{\sqrt{n!}}\left(\sqrt{\frac{m\omega}{2\hbar}}\right)^n\left(x-\frac{\hbar}{m\omega}\frac{\partial}{\partial x}\right)^n\psi_0(x)\\
&=\left(\frac{m\omega}{\pi\hbar}\right)^{1/4}\frac{1}{\sqrt{2^n n!}}\,H_n(\xi)\,e^{-\xi^2/2}
\end{align*}
其中 $\xi=x/x_0$，$H_n(\xi)=(-1)^n e^{\xi^2}\left(\dfrac{d}{d\xi}\right)^n e^{-\xi^2}$ 为 Hermite 多项式。

\subsection*{附录：有用对易关系公式}
\addcontentsline{toc}{subsection}{附录：有用对易关系公式}

对于算符函数 $\hat{f}(\hat{a},\hat{a}^\dagger)$，有：
\begin{align*}
[\hat{a},\hat{f}(\hat{a},\hat{a}^\dagger)]&=\frac{\partial\hat{f}}{\partial\hat{a}^\dagger}\\
[\hat{a}^\dagger,\hat{f}(\hat{a},\hat{a}^\dagger)]&=-\frac{\partial\hat{f}}{\partial\hat{a}}
\end{align*}
（假设 $\hat{f}$ 可展开为幂级数，且偏导表示形式导数。）

\section{离散态的表象变换}

设力学量 $\hat{A}$ 与 $\hat{B}$ 的本征方程分别为
\[
\hat{A}\ket{a_i}=a_i\ket{a_i},\qquad \hat{B}\ket{b_i}=b_i\ket{b_i},
\]
对应的本征矢各自满足完备性与正交归一条件：
\[
\sum_i\ket{a_i}\bra{a_i}=1,\qquad \sum_i\ket{b_i}\bra{b_i}=1,
\]
\[
\braket{a_i|a_j}=\delta_{ij},\qquad \braket{b_i|b_j}=\delta_{ij}.
\]

\subsection{变换算符}

定义表象变换算符
\[
\hat{S}=\sum_k\ket{a_k}\bra{b_k},
\]
则
\begin{align*}
\hat{S}\ket{b_j}&=\sum_k\ket{a_k}\braket{b_k|b_j}=\sum_k\ket{a_k}\delta_{kj}=\ket{a_j},\\[4pt]
\bra{a_j}\hat{S}&=\sum_k\braket{a_j|a_k}\bra{b_k}=\sum_k\delta_{jk}\bra{b_k}=\bra{b_j}.
\end{align*}

\subsection{态矢在不同表象中的表示}

任意态矢 $\ket{\psi}$ 在 $A$、$B$ 两表象中的列向量分别为
\[
\vec{\psi}_A=\begin{pmatrix}\braket{a_1|\psi}\\ \braket{a_2|\psi}\\ \vdots\end{pmatrix},
\qquad
\vec{\psi}_B=\begin{pmatrix}\braket{b_1|\psi}\\ \braket{b_2|\psi}\\ \vdots\end{pmatrix}.
\]

\subsection{变换矩阵的矩阵元}

在 $A$ 表象中，$\hat{S}$ 的矩阵元为
\begin{align*}
S_{ij}^{(A)}&=\bra{a_i}\hat{S}\ket{a_j}
=\sum_k\braket{a_i|a_k}\braket{b_k|a_j}
=\braket{b_i|a_j}.
\end{align*}

在 $B$ 表象中，$\hat{S}$ 的矩阵元为
\begin{align*}
S_{ij}^{(B)}&=\bra{b_i}\hat{S}\ket{b_j}
=\sum_k\braket{b_i|a_k}\braket{b_k|b_j}
=\braket{b_i|a_j}.
\end{align*}

\subsection{态矢的变换关系}

利用完备性条件，
\[
\braket{b_i|\psi}=\sum_k\braket{b_i|a_k}\braket{a_k|\psi}
=\sum_k S_{ik}\braket{a_k|\psi},
\]
即
\[
\boxed{\;\vec{\psi}_B=S\,\vec{\psi}_A\;}
\]

\subsection{力学量算符的变换关系}

算符 $\hat{H}$ 在 $B$ 表象中的矩阵元为
\begin{align*}
H_{ij}^{(B)}&=\bra{b_i}\hat{H}\ket{b_j}\\
&=\sum_{k,l}\braket{b_i|a_k}\bra{a_k}\hat{H}\ket{a_l}\braket{a_l|b_j}\\
&=\sum_{k,l}S_{ik}\,H_{kl}^{(A)}\,S_{lj}^{\dagger}.
\end{align*}
写成矩阵形式即
\[
\boxed{\;H^{(B)}=S\,H^{(A)}\,S^{\dagger}\;}
\]

\end{document}
